\section{Conclusion}
\label{sec:conclusion}

This study introduced MADKE-X, a governed debate framework that turns knowledge-enhanced multi-agent debate from an exchange of prompts that nobody can inspect into a process that is specified, checked and recorded at every step. A failure-aware interface layer, an evidence governor, a decision trace with citations and a bounded verification step were each defined by an equation, and the effect of each equation was measured on the unaltered benchmark files that are shipped with MADKE, the debate system on which the framework was built and against which it was tested.

On these files, the evidence governor places an answer-bearing passage among the first three passages of the listing for \VGainMin{} to \VGainMax{} percentage points more questions on the four question-answering datasets, with paired significance on each. It shortens the evidence listing of every selection call by \VTokRedMean\% relative to the default of the release, at a cost of less than a millisecond per question and of up to \VCoverLossMax{} points of coverage. Under identical scripted agents, MADKE-X passes all \VInvChecksAll{} pipeline checks. The summariser prompt holds the question in every case, where the release has it in none. Every selection reaches the agent that made it, where the release loses the \VSetBadRel{} selections that are written as a list. Every debate ends in the round of agreement, against \VStopPctRel\% for the release, and the consistency ratio equals the direct count, where the formula of the release reports \VConsGapMin{} to \VConsGapMax{} points too little. Each answer of MADKE-X carries a trace of the evidence delivered, the verdicts, the cited passages and a support status, and the verification step raises the number of supported answers without changing a supported one.

None of the four components depends on the debate design, the prompts or the backbone model, so the framework can in principle be placed around other debate systems that draw on retrieved evidence. The end-to-end accuracy of both systems is compared on the same questions with open-weight Qwen models under a paired protocol, and Table~\ref{tab:live} holds the results as the runs finish. Three extensions follow from the limitations of this study: a learned attribution model in place of the lexical support check, iterative retrieval for the multi-hop questions whose pools contain no answer-bearing passage, and a comparison with further model families on more questions.
